\subsection{Теоремы Вейерштрасса о непрерывной на отрезке функции.}

\textbf{Определение.} $f(x)$ непрерывна на промежутке $\langle a, b \rangle$, если $f(x)$ непрерывна в каждой точке этого промежутка.

\subsubsection*{Первая теорема Вейерштрасса (Об ограниченности)}
\textbf{Теорема.} Если функция непрерывна на отрезке, то она ограничена на нём.
\[ f(x) \in C([a,b]) \Rightarrow f(x) \in B([a,b]) \]

\begin{tcolorbox}[title=Доказательство (от противного)]
	Пусть $f(x)$ неограничена на $[a, b]$.
	Это значит, что $\forall M > 0 \; \exists x_M \in [a,b]: |f(x_M)| > M$.
	
	Построим последовательность:
	\[ \forall n \in \mathbb{N} \; \exists x_n \in [a, b]:\; |f(x_n)| > n \]
	
	Последовательность $\{x_n\} \subset [a, b]$ является ограниченной. 
	По \textbf{теореме Больцано-Вейерштрасса} из неё можно выделить сходящуюся подпоследовательность:
	\[ \exists x_{n_k} \xrightarrow[k\to \infty]{} x_0 \in [a, b] \]
	
	В силу непрерывности $f(x)$ в точке $x_0$ (определение по Гейне):
	\[ f(x_{n_k}) \to f(x_0) \]
	Следовательно, предел должен быть конечным числом.
	
	С другой стороны, из построения:
	\[ |f(x_{n_k})| > n_k \xrightarrow{k\to\infty} +\infty \]
	Получили, что $|f(x_{n_k})| \to \infty$.
	
	\textbf{Противоречие:} последовательность не может одновременно стремиться к конечному числу $f(x_0)$ и к бесконечности.
\end{tcolorbox}

\subsubsection*{Вторая теорема Вейерштрасса (О достижении точных граней)}
\textbf{Теорема.} Если $f(x) \in C([a, b])$, то $f(x)$ достигает на $[a, b]$ своих точных граней ($\max$ и $\min$).
\[ \exists x_{max}, x_{min} \in [a,b]: f(x_{max}) = \sup f(x), \; f(x_{min}) = \inf f(x) \]

\begin{tcolorbox}[title=Доказательство (для максимума)]
	1. По I теореме Вейерштрасса $f(x)$ ограничена, значит существует конечный супремум:
	\[ M = \sup_{x\in[a, b]} f(x) \in \mathbb{R} \]
	
	2. Предположим (от противного), что супремум \textbf{не достигается}.
	\[ \forall x \in [a, b]\;\; f(x) < M \]
	
	3. По определению супремума, для любого $\varepsilon > 0$ найдется $x$, где значение близко к $M$. Возьмем $\varepsilon = 1/n$:
	\[ \forall n \in \mathbb{N}\; \exists x_n \in [a, b]: M - \frac{1}{n} < f(x_n) < M \]
	
	4. Последовательность $\{x_n\} \subset [a, b]$ ограничена. По теореме Б-В выделим сходящуюся подпоследовательность:
	\[ \exists x_{n_k} \xrightarrow[k\to \infty]{} x_0 \in [a, b] \]
	
	5. Оценим предел $f(x_{n_k})$. Для всех $k$:
	\[ M - \frac{1}{n_k} < f(x_{n_k}) < M \]
	Перейдем к пределу при $k \to \infty$. По теореме о зажатой последовательности:
	\[ \lim_{k\to\infty} f(x_{n_k}) = M \]
	
	6. В силу непрерывности функции в точке $x_0$:
	\[ \lim_{k\to\infty} f(x_{n_k}) = f(x_0) \]
	
	7. \textbf{Вывод:} $f(x_0) = M$.
	Мы получили, что существует точка $x_0$, в которой функция равна $M$. Это \textbf{противоречит} нашему предположению (п. 2), что супремум не достигается.
\end{tcolorbox}

\begin{tcolorbox}[title={Ответ на вопрос: Что ломается на интервале $(a, b)$?}]
	При доказательстве мы используем теорему Больцано-Вейерштрасса, которая гарантирует, что $x_{n_k} \to x_0$.
	Для замкнутого отрезка $[a, b]$ предельная точка $x_0$ всегда принадлежит отрезку.
	
	Если же мы берем интервал $(a, b)$, то точка $x_0$ может оказаться равной $a$ или $b$ (выйти за пределы множества).
	В этом случае функция $f(x)$ не обязана быть определена или непрерывна в точке $x_0$, поэтому равенство $\lim f(x_{n_k}) = f(x_0)$ записать нельзя.
	
	\textit{Пример:} $f(x) = \frac{1}{x}$ на $(0, 1)$. Она непрерывна, но неограничена (не работает I теорема) и не достигает максимума (не работает II теорема), так как "проблемная" точка $0$ не входит в интервал.
\end{tcolorbox}
